% संपूर्ण मराठी ग्रंथातील प्रकरण 14: प्रथम-क्रम तर्कशास्त्राचा परिचय
% हा संपादनयोग्य स्रोतखंड आहे. संकलनासाठी संपूर्ण स्रोतसंग्रहातील
% अक्षररूपे, पूर्वभाग, चिन्हव्याख्या आणि आकृती-स्रोत आवश्यक आहेत.
% मूळ: Open Logic Project; मजकूर आणि रूपांतर CC BY 4.0.
% यंत्रानुवाद, दुरुस्ती आणि तपासणी: OpenAI Codex —
% GPT-5.6 Sol आणि GPT-6 Sol, दोन्ही Ultra effort.
% दुरुस्ती-पुनरावलोकन: GPT-6.1 Sol, Ultra effort.
\chapter{प्रथम-क्रम तर्कशास्त्राचा परिचय}\label{fol:int::chap}









\section{प्रथम-क्रम तर्कशास्त्र}\label{fol:int:fol:sec}

आकारिक तर्कशास्त्राशी तुमची पहिली ओळख झाली तेव्हापासून प्रथम-क्रम
तर्कशास्त्र बहुधा तुम्हाला परिचित असेल.\footnote{खरे तर, तुम्ही त्याच्याशी
परिचित आहात असे आम्ही जवळजवळ गृहीतच धरतो!{} तसे नसल्यास,
\emph{forall x} (Magnus आणि सहकारी, 2021) यासारख्या अधिक प्राथमिक
पाठ्यपुस्तकाची उजळणी करू शकता.} ते तुम्हाला ``संख्यापक तर्कशास्त्र''
किंवा ``विधेय तर्कशास्त्र'' या नावांनी माहीत असू शकते. सर्वप्रथम,
प्रथम-क्रम तर्कशास्त्र ही एक आकारिक भाषा आहे. म्हणजे तिचा एक विशिष्ट
शब्दसंग्रह असतो आणि तिच्या अभिव्यक्ती या त्या शब्दसंग्रहातील चिन्हमाला
असतात. पण प्रत्येक चिन्हमाला अनुमत नसते. अनुमत अभिव्यक्तींचे विविध
प्रकार आहेत: पदे, सूत्रे आणि वाक्ये. आपल्याला मुख्यतः
प्रथम-क्रम तर्कशास्त्रातील वाक्ये मध्ये रस आहे: ती आपल्याला
इंग्रजीतील वाक्यांचे आकारिक प्रतिरूप देतात आणि त्यांच्याविषयी
तर्कशास्त्रज्ञाला नेहमी रस असणारे प्रश्न आपण विचारू शकतो. उदाहरणार्थ:
\begin{itemize}
    \item $\psi$ हे~$\varphi$ पासून तार्किकरीत्या निष्पन्न होते का?
    \item $\varphi$ तार्किकरीत्या सत्य, तार्किकरीत्या असत्य, की
परिस्थितीवश आहे?
    \item $\varphi$ आणि $\psi$ सममूल्य आहेत का?
\end{itemize}

हे प्रश्न मुख्यतः प्रथम-क्रम तर्कशास्त्रातील वाक्यांच्या
``अर्था''विषयीचे आहेत. उदाहरणार्थ, $\psi$ हे~$\varphi$ पासून तार्किकरीत्या
निष्पन्न होते का या प्रश्नाचे तत्त्वज्ञ पुढीलप्रमाणे विश्लेषण करेल:
$\varphi$ सत्य पण~$\psi$ असत्य असणारा काही प्रसंग आहे का ($\psi$ हे~$\varphi$
पासून निष्पन्न होत नाही), की $\varphi$ सत्य करणारा प्रत्येक प्रसंग~$\psi$ ही
सत्य करतो ($\psi$ हे~$\varphi$ पासून निष्पन्न होते)? पण ``प्रसंग'' म्हणजे काय
हे अजून सांगितलेले नाही---ते \emph{चिन्हार्थमीमांसेचे} काम आहे.
प्रथम-क्रम तर्कशास्त्राची चिन्हार्थमीमांसा तत्त्वज्ञाच्या ``प्रसंग'' या
अंतःप्रज्ञ कल्पनेचे गणितीयदृष्ट्या काटेकोर प्रतिमान देते आणि---हे
महत्त्वाचे आहे---एखाद्या प्रसंगात वाक्य~$\varphi$ \emph{सत्य असणे}
म्हणजे काय याचेही प्रतिमान देते. आपण विकसित करणार असलेल्या या
गणितीयदृष्ट्या काटेकोर प्रतिमानाला संरचना म्हणतो. ``यात सत्य''
याला काटेकोर अर्थ देणाऱ्या संबंधाला \emph{पूर्ति}-संबंध म्हणतात. म्हणून
वाक्ये~$\varphi$ आणि संरचना~$\Struct{M}$ यांच्यासाठी आपण
``$\varphi$ ची~$\Struct{M}$ मध्ये पूर्ति होते'' (चिन्हांत:
$\Sat{M}{\varphi}$) हे परिभाषित करणार आहोत. हे झाल्यावर ``पासून निष्पन्न
होते'' किंवा ``तार्किकरीत्या सत्य आहे'' यांसारख्या इतर चिन्हार्थविषयक
संज्ञांच्याही काटेकोर व्याख्या देता येतील. उदाहरणार्थ,
$\lforall[x][(\varphi(x) \lif \psi(x))],
\lexists[x][\varphi(x)] \Entails \lexists[x][\psi(x)]$ आहे की नाही हे पुन्हा
गणितीय काटेकोरपणे ठरवणे या व्याख्यांमुळे शक्य होईल. अर्थातच उत्तर
``होय'' असेल. इंग्रजी वाक्यांचे प्रथम-क्रम तर्कशास्त्रात चिन्हांकन
करण्याचे प्रशिक्षण तुम्ही आधीच घेतले असेल, तर ही, उदाहरणार्थ, ``सर्व
मुंग्या कीटक आहेत, मुंग्या अस्तित्वात आहेत, म्हणून कीटक अस्तित्वात
आहेत'' यांची चिन्हांकने आहेत हे तुम्ही ओळखाल. हा उघडपणे वैध युक्तिवाद
आहे आणि म्हणून आपल्या आकारिक भाषेतील ``पासून निष्पन्न होते'' याच्या
गणितीय प्रतिमानानेही तेच उत्तर दिले पाहिजे.


आकारिक तर्कशास्त्राशी तुमची पहिली ओळख झाली तेव्हापासून आठवत असणारा
आणखी एक विषय म्हणजे \emph{निष्पत्त्या}. तुम्ही आकारिक
तर्कशास्त्राचा पहिला अभ्यासक्रम घेतला असेल, तर तुमच्या शिक्षकांनी अशा
निष्पत्त्या शोधण्याचा सराव तुमच्याकडून करून घेतला असेल; कदाचित
वरील तार्किक निष्पन्नता लागू असल्याचे दाखवणारी निष्पत्ती देखील.
निष्पत्त्या देण्याच्या अनेक पद्धती आहेत: तुम्ही ``नैसर्गिक निगमन''
किंवा ``सत्यता-वृक्ष'' नावाची पद्धत केली असेल; पण इतर अनेक पद्धतीही
आहेत. वरील तर्कशास्त्रीय प्रश्नांची उत्तरे देता येतील अशी साधने पुरवणे
हे निष्पत्ती-पद्धतींचे उद्दिष्ट आहे. उदाहरणार्थ, ज्या नैसर्गिक
निगमनातील निष्पत्तीमध्ये $\lforall[x][(\varphi(x) \lif \psi(x))]$ आणि
$\lexists[x][\varphi(x)]$ ही आधारविधाने आणि $\lexists[x][\psi(x)]$ हा
निष्कर्ष (शेवटची ओळ) आहे, ती $\lexists[x][\psi(x)]$ हे
$\lforall[x][(\varphi(x) \lif \psi(x))]$ आणि $\lexists[x][\varphi(x)]$ यांपासून
तार्किकरीत्या निष्पन्न होते याची \emph{पुष्टी करते}.

पण असे का? वरवर पाहता, निष्पत्ती-पद्धतींचा चिन्हार्थमीमांसेशी
काहीच संबंध नाही: आकारिक निष्पत्ती देणे म्हणजे केवळ काही
नियमांनी नियंत्रित रीतीने चिन्हांची मांडणी करणे; त्यात ``प्रसंगां''चा
किंवा ``यात सत्य''चा मुळीच उल्लेख नसतो. निष्पत्ती-पद्धती आणि
चिन्हार्थमीमांसा यांच्यातील संबंध अधितार्किक अन्वेषणाने प्रस्थापित
करावा लागतो. उदाहरणार्थ, आधारविधाने
$\lforall[x][(\varphi(x) \lif \psi(x))]$ आणि $\lexists[x][\varphi(x)]$ यांपासून
$\lexists[x][\psi(x)]$ ची आकारिक निष्पत्ती शक्य असते तेव्हा आणि
केवळ तेव्हाच $\lforall[x][(\varphi(x) \lif \psi(x))]$ आणि
$\lexists[x][\varphi(x)]$ ही दोन्ही मिळून
$\lexists[x][\psi(x)]$ तार्किकरीत्या निष्पन्न करतात, याची गणितीय
सिद्धता आवश्यक आहे. मात्र, हे करता येण्यापूर्वी विन्यासमीमांसा आणि
चिन्हार्थमीमांसा यांच्या व्याख्या अचूक करण्याचे बरेच कष्टप्रद काम
करावे लागते.












\section{विन्यासमीमांसा}\label{fol:int:syn:sec}

प्रथम-क्रम तर्कशास्त्रातील कोणत्या चिन्हमाला वाक्ये म्हणून गणल्या
जातात हे प्रथम काटेकोरपणे सांगावे लागेल. हे आपण नंतर करू; सध्या
उदाहरणांच्या साहाय्याने पुढे जाऊ. प्रथम-क्रम तर्कशास्त्राचे मूलभूत
रचनाघटक---म्हणजे त्याचा शब्दसंग्रह---दोन भागांत विभागलेले आहेत. पहिल्या
भागात विशिष्ट गोष्टींबद्दल काही सांगण्यासाठी किंवा विशिष्ट गोष्टी
निर्देशित करण्यासाठी आपण वापरतो ती चिन्हे येतात. गोष्टी निर्देशित
करण्यासाठी आपण स्थिरांक वापरतो आणि निर्देशित केलेल्या गोष्टींबद्दल
काही सांगण्यासाठी विधेये वापरतो. उदाहरणार्थ, एखादी एक गोष्ट
निर्देशित करण्यासाठी $\Obj a$ हे स्थिरांक म्हणून वापरू शकतो आणि मग
वाक्य~$\Atom{\Obj P}{\Obj a}$ वापरून तिच्याबद्दल काही सांगू शकतो.
``$\Obj a$'' आणि ``$\Obj P$'' यांचे काही अर्थ तुम्ही मनात धरले असतील,
तर $\Atom{\Obj P}{\Obj a}$ हे इंग्रजीतील वाक्य म्हणून वाचू शकता (आणि
आकारिक तर्कशास्त्र प्रथम शिकताना तुम्ही बहुधा तसे केलेही असेल).
प्रथम-क्रम तर्कशास्त्रातील अशी साधी वाक्ये मिळाल्यानंतर,
शब्दसंग्रहातील दुसरा भाग---तार्किक चिन्हे (संयोजक आणि संख्यापक)---वापरून
त्यांपासून अधिक जटिल वाक्ये रचता येतात. म्हणून, उदाहरणार्थ,
$(\Atom{\Obj P}{\Obj a} \land \Atom{\Obj Q}{\Obj b})$ किंवा
~$\lexists[\Obj x][\Atom{\Obj P}{\Obj x}]$ यांसारख्या अभिव्यक्ती
रचता येतात.

काटेकोर अधितार्किक अभ्यासासाठी आवश्यक असलेल्या चिन्हार्थमीमांसेच्या
नेमक्या व्याख्या आणि आपल्या निष्पत्ती-पद्धतींचे नियम देता यावेत,
यासाठी प्रथम-क्रम तर्कशास्त्रातील वाक्य म्हणून नेमके काय गणले
जाते याची काटेकोर व्याख्या सर्वप्रथम द्यावी लागेल. मूलभूत कल्पना समजणे
पुरेसे सोपे आहे: फक्त विधेये आणि स्थिरांक यांपासून आपण
$\Atom{\Obj P}{\Obj a}$ यासारखी काही साधी वाक्ये रचू शकतो. मग
संयोजक आणि संख्यापक वापरून त्यांपासून अधिक जटिल वाक्ये रचतो. पण अधिक
जटिल वाक्ये रचण्याची अनुमती देणारे नियम नेमके कोणते? ते स्पष्ट
करावेच लागतील; अन्यथा ``प्रथम-क्रम तर्कशास्त्रातील वाक्य'' याची
आपण पुरेशी काटेकोर व्याख्या केलेली नाही. येथे काही अडचणी आहेत. पहिली
अडचण म्हणजे योग्य चिन्हमालाच वाक्ये म्हणून गणल्या जातील याची
खात्री करणे. दुसरी म्हणजे हे अशा रीतीने करणे की
\emph{सर्व}~वाक्ये विषयी गणितीय सिद्धता देता येतील. शेवटी,
``$\varphi$ मधील प्रत्येक $\Obj x$ च्या जागी~$\Obj b$ ठेवा'' यासारख्या
वाक्ये वरील काही प्राथमिक क्रियांच्याही काटेकोर व्याख्या द्याव्या
लागतील. प्रथम-क्रम तर्कशास्त्रातील संख्यापक आणि चले यांमुळे हे
कसे करावे हे पूर्णपणे स्पष्ट नसते, ही अडचण आहे. उदाहरणार्थ,
$\lexists[\Obj x][\Atom{\Obj P}{\Obj a}]$ हे वाक्य म्हणून
गणले जावे का? $\lexists[\Obj x][\lexists[\Obj x][\Atom{\Obj P}{\Obj
x}]]$ चे काय? आणि ``$(\Atom{\Obj P}{\Obj x} \land \lexists[\Obj
x][\Atom{\Obj P}{\Obj x}])$ मध्ये $\Obj x$ च्या जागी~$\Obj b$ ठेवा''
याचा परिणाम काय असावा?












\section{सूत्र}\label{fol:int:fml:sec}

प्रथम-क्रम तर्कशास्त्रातील वाक्ये काटेकोरपणे निर्दिष्ट करण्यासाठी
आणि चलांच्या वापरामुळे निर्माण होणाऱ्या अडचणी हाताळण्यासाठी
आपण पुढील पद्धत वापरणार आहोत. प्रथम आपण अभिव्यक्तींचा एक
\emph{वेगळा} संच परिभाषित करतो: सूत्रे. ते केल्यानंतर त्यांत
चले कोणती भूमिका बजावतात याचा विचार करता येतो---आणि
``मुक्त'' व ``बद्ध'' चले या काही इतर कल्पनांच्या आधारे
वाक्य म्हणजे काय हे परिभाषित करता येते (म्हणजे, मुक्त
चले नसलेले सूत्र). यामुळे ``वाक्य'' ची व्याख्या
अधिक हाताळण्याजोगी होते एवढ्याचसाठी आपण हे करत नाही; पूर्ति ही
चिन्हार्थविषयक संकल्पना ज्या रीतीने परिभाषित करू त्यासाठीही ते
निर्णायक ठरेल.

फक्त एक विधेय~$\Obj P$, एक स्थिरांक~$\Obj a$ आणि
$\lnot$, $\land$ व~$\lexists$ ही तार्किक चिन्हे असलेल्या एका साध्या
प्रथम-क्रम भाषेसाठी ``सूत्र'' परिभाषित करू या. आपल्या संपूर्ण
व्याख्या यापेक्षा खूप व्यापक असतील: आपण अनंतपणे अनेक विधेये
आणि स्थिरांक स्वीकारू. किंबहुना, स्थिरांक आणि चले
यांच्याशी जोडून ``पदे'' रचता येतील अशा फलनांचाही विचार करू.
सध्या $\Obj a$ आणि चरे एवढीच आपली पदे असतील. आपल्याला अनंतपणे अनेक
चले मात्र आवश्यक आहेत. अधिकृतपणे आपण $\Obj v_0$, $\Obj
v_1$, \dots, ही चिन्हे चरे म्हणून वापरू.

\begin{defn}
\emph{सूत्रे} चा संच~$\Frm$ पुढीलप्रमाणे परिभाषित केला आहे:
\begin{enumerate}
\item\label{fol:int:fml:fmls-atom} $\Atom{\Obj P}{\Obj a}$ आणि $\Atom{\Obj
  P}{\Obj v_i}$ ही सूत्रे आहेत ($i \in \Nat$).

\item \label{fol:int:fml:fmls-not}$\varphi$ हे सूत्र असेल, तर $\lnot \varphi$ हे
  सूत्र आहे.

\item $\varphi$ आणि $\psi$ ही सूत्रे असतील, तर $(\varphi \land
  \psi)$ हे सूत्र आहे.

\item \label{fol:int:fml:fmls-ex}$\varphi$ हे सूत्र आणि $x$ हे चल
  असेल, तर $\lexists[x][\varphi]$ हे सूत्र आहे.

\item \label{fol:int:fml:fmls-limit}दुसरे काहीही सूत्र नाही.
\end{enumerate}
\end{defn}

कोणत्याही $i \in \Nat$ साठी $\Atom{\Obj P}{\Obj a}$ आणि
$\Atom{\Obj P}{\Obj v_i}$ ही सूत्रे आहेत, असे
\ref{fol:int:fml:fmls-atom} सांगतो. त्यांना तथाकथित \emph{आण्विक} सूत्रे
म्हणतात. त्यांपासून सुरुवात करता येते. इतर उपवाक्ये आपण आधीच
रचलेल्यांपासून नवी सूत्रे रचण्याचे मार्ग देतात. म्हणून,
उदाहरणार्थ, $\Atom{\Obj P}{\Obj v_2}$ हे \ref{fol:int:fml:fmls-atom} नुसार आधीच
सूत्र असल्यामुळे, $\lnot \Atom{\Obj P}{\Obj v_2}$ हे
\ref{fol:int:fml:fmls-not} नुसार सूत्र आहे. मग
\ref{fol:int:fml:fmls-ex} नुसार $\lexists[\Obj v_2][\lnot \Atom{\Obj P}{\Obj
v_2}]$ हे आणखी एक सूत्र आहे; आणि असेच पुढे जाता येते. \emph{फक्त} याच
रीतीने रचता येणाऱ्या चिन्हमाला सूत्रे म्हणून गणल्या जातात, असे
\ref{fol:int:fml:fmls-limit} सांगतो. विशेषतः,
$\lexists[\Obj v_0][\Atom{\Obj P}{\Obj a}]$ आणि $\lexists[\Obj
  v_0][\lexists[\Obj v_0][\Atom{\Obj P}{\Obj a}]]$ ही
\emph{खरोखर} सूत्रे म्हणून गणली जातात; परंतु अनावश्यक बाहेरील
कंसांमुळे $(\lnot \Atom{\Obj P}{\Obj a})$ हे गणले जात नाही.

सूत्रे परिभाषित करण्याच्या या पद्धतीला \emph{विगमनाधारित
व्याख्या} म्हणतात आणि त्यामुळे \emph{संरचनात्मक विगमन} म्हणतात त्या
विगमनाने सिद्धता करण्याच्या प्रकाराचा वापर करून सूत्रे विषयी
गोष्टी सिद्ध करता येतात. यांची सर्वसाधारण चर्चा
\ref{mth:ind:idf:sec} आणि \ref{mth:ind:sti:sec} येथे केली
आहे; पुढील सिद्धतांचा सखोल अभ्यास करण्यापूर्वी तिची उजळणी करावी. मूलतः,
सर्व सूत्रे साठी एखादी गोष्ट सत्य आहे याची सिद्धता द्यायची असेल,
तर ती आण्विक सूत्रे साठी सत्य आहे हे प्रथम दाखवता आणि नंतर ती
कोणत्याही सूत्र~$\varphi$ साठी (आणि~$\psi$ साठी) सत्य \emph{असेल}, तर
ती $\lnot \varphi$, $(\varphi \land \psi)$ आणि $\lexists[x][\varphi]$ यांच्यासाठीही
\emph{सत्य असते}, हे दाखवता. उदाहरणार्थ, यावरून ती
$\lexists[\Obj v_2][\lnot \Atom{\Obj P}{\Obj v_2}]$ साठी सत्य आहे हे
सिद्ध होते: पहिल्या भागावरून ती आण्विक सूत्र~$\Atom{\Obj
P}{\Obj v_2}$ साठी सत्य आहे हे माहीत असते. मग दुसऱ्या भागावरून ती
$\lnot \Atom{\Obj P}{\Obj v_2}$ साठी सत्य आहे आणि त्याच भागावरून पुन्हा
$\lexists[\Obj v_2][\lnot \Atom{\Obj P}{\Obj v_2}]$ याच्यासाठी सत्य
आहे, हे मिळते. सर्व सूत्रे आण्विक सूत्रांपासून
विगमनाधारित रीतीने निर्माण केली जात असल्यामुळे, ही पद्धत त्यांपैकी
कोणत्याही सूत्रासाठी चालते.












\section{पूर्ति}\label{fol:int:sat:sec}

सूत्रे म्हणजे काय हे माहीत झाल्यावर आपण लगेच प्रथम-क्रम
तर्कशास्त्राच्या चिन्हार्थमीमांसेकडे वळू शकतो: येथे मूलभूत व्याख्या
संरचनेची आहे. आपल्या साध्या भाषेसाठी संरचना~$\Struct M$
चे फक्त तीन घटक असतात: \emph{प्रांत} म्हणतात तो अरिक्त संच
$\Domain{M}$, $\Struct M$ मध्ये $\Obj a$ ने निर्देशित केलेली गोष्ट, आणि
$\Struct M$ मध्ये $\Obj P$ ज्या गोष्टींबाबत सत्य आहे त्या गोष्टी.
$\Obj a$ ने निर्देशित केलेली वस्तू~$\Assign{\Obj a}{M}$ ने, तर
$\Obj P$ ज्या गोष्टींबाबत सत्य आहे त्यांचा संच~$\Assign{\Obj P}{M}$ ने
दर्शवतात. संरचना~$\Struct{M}$ मध्ये याच तीन गोष्टी असतात:
$\Domain{M}$, $\Assign{\Obj a}{M} \in \Domain{M}$ आणि
$\Assign{\Obj P}{M} \subseteq \Domain{M}$. सर्वसाधारण स्थिती अधिक
गुंतागुंतीची असेल, कारण त्यात अनेक विधेये आणि स्थिरांक
असतील, विधेये ना एकाहून अधिक स्थाने असू शकतील, आणि
फलने ही असतील.


ज्या सूत्रे मध्ये चले नाहीत त्यांच्यासाठी पूर्तिची
व्याख्या देण्यास हे पुरेसे आहे. आपण सूत्रे ज्या रीतीने परिभाषित
केली तिचे प्रतिबिंब दाखवणारी विगमनाधारित व्याख्या देण्याची कल्पना आहे.
एखाद्या आण्विक सूत्राची~$\Struct{M}$ मध्ये पूर्ति केव्हा होते हे आपण
सांगतो आणि मग, उदाहरणार्थ, $\varphi$ ची~$\Struct{M}$ मध्ये पूर्ति होते की
नाही याच्या आधारे $\lnot \varphi$ ची~$\Struct{M}$ मध्ये पूर्ति केव्हा होते
हे सांगतो. उदाहरणार्थ, पुढीलप्रमाणे व्याख्या करता येईल:
\begin{enumerate}
  \item $\Atom{\Obj P}{\Obj a}$ ची~$\Struct{M}$ मध्ये पूर्ति तेव्हा आणि
  केवळ तेव्हाच होते, जेव्हा $\Assign{\Obj a}{M} \in \Assign{\Obj P}{M}$.
  \item $\lnot \varphi$ ची~$\Struct{M}$ मध्ये पूर्ति तेव्हा आणि केवळ तेव्हाच
  होते, जेव्हा $\varphi$ ची~$\Struct{M}$ मध्ये पूर्ति होत नाही.
  \item $(\varphi \land \psi)$ ची~$\Struct{M}$ मध्ये पूर्ति तेव्हा आणि केवळ
  तेव्हाच होते, जेव्हा $\varphi$ ची~$\Struct{M}$ मध्ये पूर्ति होते आणि
  $\psi$ ची~$\Struct{M}$ मध्येही पूर्ति होते.
\end{enumerate}
$\Domain{M} = \{0, 1, 2\}$, $\Assign{\Obj a}{M} = 1$ आणि
$\Assign{\Obj P}{M} = \{1, 2\}$ असे समजू. ही व्याख्या आपल्याला
$\Atom{\Obj P}{\Obj a}$ ची~$\Struct{M}$ मध्ये पूर्ति होते असे सांगेल
(कारण $\Assign{\Obj a}{M} =  1 \in \{1,2\} = \Assign{\Obj P}{M}$).
ती पुढे असे सांगते की $\lnot \Atom{\Obj P}{\Obj a}$ ची~$\Struct{M}$
मध्ये पूर्ति होत नाही; म्हणून $\lnot\lnot \Atom{\Obj P}{\Obj a}$ ची
पूर्ति होते आणि $(\lnot \Atom{\Obj P}{\Obj a} \land \Atom{\Obj P}{\Obj a})$
ची पूर्ति होत नाही; आणि असेच पुढे जाता येते.

संख्यापकांसाठी व्याख्या द्यायची झाल्यावर अडचण निर्माण होते: आपल्याला
``$\lexists[\Obj v_0][\Atom{\Obj P}{\Obj v_0}]$ ची पूर्ति तेव्हा आणि
केवळ तेव्हाच होते, जेव्हा $\Atom{\Obj P}{\Obj v_0}$ ची पूर्ति होते''
असे काहीतरी म्हणायचे आहे. पण संरचना~$\Struct{M}$ हे
चले विषयी काय करायचे ते सांगत नाही. आपल्याला प्रत्यक्षात
$\Atom{\Obj P}{\Obj v_0}$ ची \emph{$\Obj v_0$ च्या एखाद्या मूल्यासाठी}
पूर्ति होते असे म्हणायचे आहे. हे काटेकोर करण्यासाठी फक्त $\Obj a$ लाच
नव्हे, तर~$\Obj v_0$ लाही~$\Domain{M}$ मधील घटक नेमून देण्याचा
मार्ग हवा. यासाठी आपण चल ची \emph{मूल्यांकने} आणतो.
चल चे मूल्यांकन म्हणजे चले ना~$\Domain{M}$ मधील
घटक कडे पाठवणारे फलन~$s$ (आपल्या उदाहरणात, $0$, $1$ किंवा~$2$
यांपैकी एका घटकाकडे). एखाद्या सूत्रामध्ये कोणती चले
येतील हे आपल्याला आधी माहीत नसल्यामुळे, $s$ कोणत्या चले ना
मूल्ये नेमून देते यावर मर्यादा घालता येत नाही. सोपा उपाय म्हणजे $s$ ने
\emph{सर्व} चले~$\Obj v_0$, $\Obj v_1$, \dots\@ यांना मूल्ये
नेमून द्यावीत अशी अट ठेवणे. आपल्याला लागतील तीच मूल्ये आपण वापरू.


सूत्रांची पूर्ति फक्त संरचनेच्या सापेक्ष परिभाषित
करण्याऐवजी, आपण ती संरचना~$\Struct{M}$ \emph{आणि}
चल चे मूल्यांकन~$s$ या दोन्हींच्या सापेक्ष परिभाषित करू आणि
संक्षेपाने $\Sat{M}{\varphi}[s]$ असे लिहू. चले असलेली आण्विक
सूत्रे हाताळण्यासाठी आपल्या व्याख्येत आता एक अतिरिक्त उपवाक्य
असेल:
\begin{enumerate}
  \item $\Sat{M}{\Atom{\Obj P}{\Obj a}}[s]$ तेव्हा आणि केवळ तेव्हाच,
  जेव्हा $\Assign{\Obj a}{M} \in \Assign{\Obj P}{M}$.
  \item $\Sat{M}{\Atom{\Obj P}{\Obj v_i}}[s]$ तेव्हा आणि केवळ तेव्हाच,
  जेव्हा $s(\Obj v_i) \in \Assign{\Obj P}{M}$.
  \item $\Sat{M}{\lnot \varphi}[s]$ तेव्हा आणि केवळ तेव्हाच, जेव्हा
  $\Sat{M}{\varphi}[s]$ नाही.
  \item $\Sat{M}{(\varphi \land \psi)}[s]$ तेव्हा आणि केवळ तेव्हाच, जेव्हा
  $\Sat{M}{\varphi}[s]$ आणि $\Sat{M}{\psi}[s]$ दोन्ही आहेत.
\end{enumerate}
ठीक आहे, यामुळे एक अडचण सुटते: $s(\Obj v_0)$ या मूल्यासाठी
$\Struct{M}$ हे $\Atom{\Obj P}{\Obj v_0}$ ची पूर्ति केव्हा करते हे आता
सांगता येते. $\lexists[\Obj v_0][\Atom{\Obj P}{\Obj v_0}]$ साठी योग्य
व्याख्या मिळवण्यासाठी आणखी एक गोष्ट करावी लागेल:
$\Sat{M}{\lexists[\Obj v_0][\Atom{\Obj P}{\Obj v_0}]}[s]$ तेव्हा आणि
केवळ तेव्हाच असावे, जेव्हा~$\Obj v_0$ ला मूल्य नेमून देणाऱ्या
\emph{एखाद्या} पद्धती~$s'$ साठी
$\Sat{M}{\Atom{\Obj P}{\Obj v_0}}[s']$ असते. पण~$\Obj v_0$ ला नेमलेले
मूल्य $s(\Obj v_0)$ ने निर्देशित केलेले मूल्यच असणे आवश्यक नाही.
त्यासाठी आपण एक संकेतपद्धती आणू: $m \in \Domain{M}$ असेल, तर
$\Subst{s}{m}{\Obj v_0}$ हे~$s$ सारखेच असणारे (फक्त~$\Obj v_0$ सोडून
इतर सर्व चले साठी) पण~$\Obj v_0$ ला~$m$ नेमून देणारे
मूल्यांकन मानू. आता आपली व्याख्या अशी देता येईल:
\begin{enumerate}\setcounter{enumi}{4}
  \item $\Sat{M}{\lexists[\Obj v_i][\varphi]}[s]$ तेव्हा आणि केवळ तेव्हाच,
  जेव्हा एखाद्या~$m \in \Domain{M}$ साठी
  $\Sat{M}{\varphi}[\Subst{s}{m}{\Obj v_i}]$.
\end{enumerate}
हे योग्य रीतीने चालते का? प्रत्येक~$i \in \Nat$ साठी
$s(\Obj v_i) = 0$ ठेवू या. एखादा $m \in \Domain{M}$ असा असेल की
$\Sat{M}{\Atom{\Obj P}{\Obj v_0}}[\Subst{s}{m}{\Obj v_0}]$, तेव्हा आणि
केवळ तेव्हाच $\Sat{M}{\lexists[\Obj v_0][\Atom{\Obj P}{\Obj v_0}]}[s]$.
आणि असा घटक आहे: आपण $m = 1$ किंवा $m = 2$ निवडू शकतो. $s$ नेच
~$\Obj v_0$ ला नेमून दिलेले $s(\Obj v_0)$ हे मूल्य---या उदाहरणात,
$0$---हे काम करत नसले तरी हे सत्य आहे, हे लक्षात घ्या. आपल्याकडे
$\Sat{M}{\Atom{\Obj P}{\Obj v_0}}[\Subst{s}{1}{\Obj v_0}]$ आहे, पण
$\Sat{M}{\Atom{\Obj P}{\Obj v_0}}[s]$ नाही.

हे गोंधळात टाकणारे आणि अवजड वाटत असेल, तर ते तसे आहे. पण सर्व
सूत्रे साठी पूर्तिची काटेकोर, विगमनाधारित व्याख्या देण्यासाठी ही
अतिरिक्त गुंतागुंत आवश्यक आहे आणि चिन्हार्थविषयक संकल्पना काटेकोरपणे
परिभाषित करण्यासाठी आपल्याला अशाच काहीतरी पद्धतीची गरज आहे. हे
करण्याचे इतरही मार्ग आहेत, पण ते सर्व तितकेच (अन)सुबक आहेत.












\section{वाक्य}\label{fol:int:snt:sec}

ठीक आहे, आता आपल्याकडे संरचना आणि सूत्रे यांच्यासाठी
पूर्तिच्या (``यात सत्य'') व्याख्येचा आराखडा आहे. पण तिला हा अतिरिक्त
भाग---चल चे मूल्यांकन---आवश्यक आहे आणि आपल्याला हवी होती ती
वाक्यांची व्याख्या. मूल्यांकने कशी काढून टाकायची आणि
वाक्ये म्हणजे काय?

आकारिक तर्कशास्त्राशी तुमची पहिली ओळख झाली तेव्हा चले आणि
संख्यापक यांच्यातील संबंधाची केलेली चर्चा तुम्हाला बहुधा आठवत असेल.
संख्यापकानंतर नेहमी चल येते आणि मग त्या संख्यापकाचा
वाक्य मधील ज्या भागाला उपयोग होतो (त्याची ``व्याप्ती''), त्या
भागात ते चल त्या संख्यापकाने ``बद्ध'' केले आहे असे आपण
समजतो. सूत्रे मध्ये प्रत्येक चल ला जुळणारा संख्यापक
असणे आवश्यक नव्हते आणि जुळणारे संख्यापक नसलेली चले ही
``मुक्त'' किंवा ``अबद्ध'' असतात. मुक्त चले नसलेली सर्व
सूत्रे म्हणजे वाक्ये असे आपण मानू.

सूत्र~$\varphi$ मधील चल ची एखादी आस्थिती केव्हा बद्ध असते,
तिला कोणता संख्यापक बद्ध करतो आणि ती केव्हा मुक्त असते, याची अंतःप्रज्ञ
कल्पना समजणे पुन्हा अवघड नाही. हे तपासण्याची, कदाचित कंस मोजण्याची,
एखादी पद्धत तुम्ही शिकला असाल. पण आपण काटेकोर व्याख्येचा आग्रह धरला
पाहिजे---आणि सूत्रांची व्याख्या आपण विगमनाने केली असल्यामुळे,
सूत्र~$\varphi$ मधील चल~$x$ च्या मुक्त आणि बद्ध आस्थितींची
व्याख्याही विगमनाने देता येते. उदाहरणार्थ, आपल्या सोप्या भाषेसाठी ती
अशी दिसू शकेल:
\begin{enumerate}
  \item $\varphi$ हे आण्विक असेल, तर त्यातील $x$ च्या सर्व आस्थिती मुक्त
  असतात (म्हणजे, $\Atom{\Obj P}{x}$ मधील $x$ ची आस्थिती मुक्त असते).
  \item $\varphi$ हे $\lnot \psi$ या रूपाचे असेल, तर $\lnot \psi$ मधील~$x$ ची
  आस्थिती तेव्हा आणि केवळ तेव्हाच मुक्त असते, जेव्हा~$\psi$ मधील~$x$ ची
  संबंधित आस्थिती मुक्त असते (म्हणजे,~$\psi$ मधील चरांच्या मुक्त
  आस्थिती म्हणजे नेमक्या $\lnot \psi$ मधील संबंधित आस्थिती होत).
  \item $\varphi$ हे $(\psi \land \chi)$ या रूपाचे असेल, तर $(\psi \land \chi)$
  मधील~$x$ ची आस्थिती तेव्हा आणि केवळ तेव्हाच मुक्त असते, जेव्हा~$\psi$
  किंवा~$\chi$ मधील~$x$ ची संबंधित आस्थिती मुक्त असते.
  \item $\varphi$ हे $\lexists[x][\psi]$ या रूपाचे असेल, तर $\varphi$ मधील $x$ ची
  कोणतीही आस्थिती मुक्त नसते; ते $\lexists[y][\psi]$ या रूपाचे असेल आणि
  $y$ हे~$x$ पेक्षा वेगळे चल असेल, तर
  $\lexists[y][\psi]$ मधील~$x$ ची आस्थिती तेव्हा आणि केवळ तेव्हाच मुक्त
  असते, जेव्हा~$\psi$ मधील~$x$ ची संबंधित आस्थिती मुक्त असते.
\end{enumerate}

चरांच्या मुक्त आणि बद्ध आस्थितींची काटेकोर व्याख्या मिळाल्यानंतर आपण
सरळ असे म्हणू शकतो: मुक्त चलांची आस्थिती नसलेले कोणतेही
सूत्र म्हणजे वाक्य होय.












\section{चिन्हार्थविषयक संकल्पना}\label{fol:int:sem:sec}

$\Sat{M}{\varphi}[s]$ लागू आहे का याचा विचार करताना, सोयीसाठी आपण $s$ कडून
सर्व चले ना मूल्ये नेमून घेतो, पण~$\varphi$ मधील चले ना
त्याने नेमून दिलेली मूल्येच फक्त वापरली जातात, असे आपण वर म्हटले.
खरे तर, $\varphi$ मधील \emph{मुक्त} चरांची मूल्येच महत्त्वाची असतात. आपण
काटेकोर असल्यामुळे अर्थातच हे तथ्य सिद्ध करणार आहोत. वाक्ये मध्ये
मुक्त चरे नसल्यामुळे, संरचनांमध्ये त्यांची पूर्ति होते की नाही
यासाठी $s$ चा काहीही फरक पडत नाही. म्हणून, $\varphi$ हे वाक्य असेल,
तेव्हा ``प्रत्येक~$s$ साठी $\Sat{M}{\varphi}[s]$'' असा $\Sat{M}{\varphi}$ चा अर्थ
परिभाषित करता येतो; आणि तो योगायोगाने किमान एका~$s$ साठी
$\Sat{M}{\varphi}[s]$ असते तेव्हा आणि केवळ तेव्हाच सत्य असतो. सूत्रे
साठी पूर्तिची कार्यक्षम व्याख्या मिळवण्यासाठी आपल्याला चल ची
मूल्यांकने आणावी लागतात, पण वाक्ये साठी पूर्ति ही चल
च्या मूल्यांकनांपासून स्वतंत्र असते.

``$\Sat{M}{\varphi}$'' ची व्याख्या मिळाल्यानंतर प्रथम-क्रम तर्कशास्त्रातील
वाक्यांच्या संदर्भात ``प्रसंग'' आणि ``यात सत्य'' यांचा अर्थ काय
हे आपल्याला माहीत होते. मग वाक्ये साठीच्या $\Sat{M}{\varphi}$ च्या
व्याख्येच्या आधारे वैधता, तार्किक निष्पन्नता आणि पूर्ततायोग्यता या मूलभूत
चिन्हार्थविषयक संकल्पना परिभाषित करता येतात. प्रत्येक संरचना
एखाद्या वाक्याची पूर्ति करत असेल, तर ते वाक्य वैध, $\Entails \varphi$, असते.
~$\Gamma$ मधील सर्व वाक्यांची पूर्ति करणारे प्रत्येक
संरचना हे~$\varphi$ चीही पूर्ति करत असेल, तर वाक्यांच्या संचातून
ते तार्किकरीत्या निष्पन्न होते, $\Gamma \Entails \varphi$. आणि एखादे
संरचना एखाद्या वाक्यांच्या संचातील सर्व वाक्यांची एकाच
वेळी पूर्ति करत असेल, तर तो संच पूर्ततायोग्य असतो.

सूत्रे विगमनाधारित रीतीने परिभाषित केली आहेत आणि पूर्ति हीही
सूत्रांच्या संरचनेवरील विगमनाने परिभाषित केली आहे; म्हणून आपल्या
चिन्हार्थमीमांसेचे गुणधर्म सिद्ध करण्यासाठी आणि परिभाषित
चिन्हार्थविषयक संकल्पनांचा परस्परसंबंध लावण्यासाठी विगमन वापरता येते.
यांपैकी काही गुणधर्म आपण एकत्र करून सिद्ध करू; काही अंशी प्रत्येक
गुणधर्म स्वतंत्रपणे रोचक असल्यामुळे, पण मुख्यतः चिन्हार्थमीमांसा आणि
निष्पत्ती-पद्धती यांच्यातील संबंधाचा अभ्यास पुढे करताना त्यांपैकी
अनेक उपयोगी पडणार असल्यामुळे. हे करण्यासाठी आपण अजून उल्लेख न केलेल्या
काही विन्यासविषयक संकल्पना आणि क्रियाही (काटेकोरपणे, म्हणजे विगमनाने)
परिभाषित कराव्या लागतील.












\section{आदेशन}\label{fol:int:sub:sec}

सर्व गोष्टींचा परस्परसंबंध कसा आहे आणि विन्यासमीमांसा व
चिन्हार्थमीमांसा यांचा विकास पुढील अधिक प्रगत अभ्यासाचा पाया कसा घालतो,
हे दाखवण्यासाठी एका उदाहरणाची चर्चा करू. आपल्या निष्पत्ती-पद्धतींनी
आपल्याला $\lforall[\Obj v_0][\Atom{\Obj P}{\Obj v_0}]$ पासून
$\Atom{\Obj P}{\Obj a}$ निष्पन्न करू दिले पाहिजे. कदाचित हे निगमन नियम
म्हणूनही मांडायचे असेल. पण त्यासाठी ते शक्य तितक्या सर्वसाधारण शब्दांत
मांडता आले पाहिजे: फक्त $\Obj P$, $\Obj a$ आणि $\Obj v_0$ यांच्यासाठी
नव्हे, तर कोणत्याही सूत्र~$\varphi$, पद~$t$ आणि चल~$x$
यांच्यासाठी. (स्थिरांक ही पदे असतात हे आठवा; पण स्थिरांक आणि
फलने यांच्यापासून रचलेल्या अधिक जटिल पदांचाही आपण विचार करू.)
म्हणून आपल्याला असे काहीतरी म्हणता आले पाहिजे: ``जेव्हा तुम्ही
$\lforall[x][\varphi(x)]$ निष्पन्न केले असेल, तेव्हा~$\varphi(t)$ चे अनुमान
करणे समर्थित असते---$\lforall[x]$ काढून आणि~$x$ च्या जागी~$t$ ठेवून
मिळणारे हे सूत्र आहे.'' पण ``$x$ च्या जागी~$t$ ठेवणे'' म्हणजे नेमके
काय? $\varphi(x)$ आणि~$\varphi(t)$ यांच्यातील संबंध कोणता? हे नेहमी चालते का?


हे काटेकोर करण्यासाठी आपण \emph{आदेशन} ही क्रिया परिभाषित करतो.
आदेशन प्रत्यक्षात गुंतागुंतीचे आहे, कारण~$\varphi$ मधील सर्व~$x$ फक्त~$t$ ने
बदलता येत नाहीत आणि प्रत्येक~$t$ चे कोणत्याही~$x$ साठी आदेशन करता येत
नाही. आपण हे पुन्हा विगमनाधारित व्याख्या वापरून हाताळू. पण हे
केल्यानंतर ``$\lforall[x][\varphi(x)]$ पासून $\varphi(t)$ चे अनुमान करा'' असा
निगमन नियम निर्दिष्ट करणे ही काटेकोर व्याख्या बनते. शिवाय,
$\lforall[x][\varphi(x)]$ पासून~$\varphi(t)$ तार्किकरीत्या निष्पन्न होते या अर्थाने
हा चांगला निगमन नियम आहे, हे दाखवता येईल. पण हे सिद्ध करण्यासाठी प्रथम
दर्शनी ``आपण हे का करत आहोत?'' असा प्रश्न पडेल अशी आणखी एक गोष्ट पुन्हा
सिद्ध करावी लागेल. $\lforall[x][\varphi(x)]$ पासून~$\varphi(t)$ तार्किकरीत्या
निष्पन्न होते हे पुढील तथ्यावर अवलंबून आहे: $\Sat{M}{\varphi(t)}$ लागू आहे
की नाही हे फक्त पद~$t$ च्या मूल्यावर अवलंबून असते. म्हणजे, $t$ ने
निर्देशित होणारा~$\Domain{M}$ मधील कोणताही घटक आपण $m$ मानला,
तर $\Sat{M}{\varphi(t)}[s]$ तेव्हा आणि केवळ तेव्हाच, जेव्हा
$\Sat{M}{\varphi(x)}[\Subst{s}{m}{x}]$. $t$ मध्ये चले असली तरीही
हे लागू असते; पण हा परिणाम नेमका कसा मांडायचा याबाबत काळजी घ्यावी
लागेल.












\section{प्रतिमाने आणि उपपत्ती}\label{fol:int:mod:sec}

प्रथम-क्रम तर्कशास्त्राची विन्यासमीमांसा आणि चिन्हार्थमीमांसा परिभाषित
केल्यानंतर संरचनेचे गुणधर्म आणि चिन्हार्थविषयक संकल्पना यांचा
अभ्यास सुरू करता येतो. निष्पत्ती-पद्धतीही परिभाषित करून त्यांचा
अभ्यास करता येतो. वाक्यांचा एखादा संच घेतल्यावर आपण विचारू
शकतो: त्या संचातील सर्व वाक्ये कोणती संरचना सत्य करतात?
वाक्यांचा संच~$\Gamma$ दिला असता, त्यांची पूर्ति करणाऱ्या
संरचना~$\Struct{M}$ ला \emph{$\Gamma$ चे प्रतिमान} म्हणतात.
~$\Gamma$ पासून सुरुवात करून आपण त्याची प्रतिमाने शोधू शकतो---ती कशी
दिसतात? ती किती मोठी किंवा लहान असावी लागतात? पण एखादे एक
संरचना किंवा संरचनेचा संग्रह घेऊनही विचारता येते: त्यांत
कोणती वाक्ये सत्य आहेत? या संरचनांमध्ये ती आणि फक्त तीच
सत्य असतात या अर्थाने त्यांचे \emph{अभिलक्षण} करणारी वाक्ये आहेत
का? अशा प्रकारचे प्रश्न \emph{प्रतिमान उपपत्ती} या क्षेत्रातील आहेत.
त्यांवरच \emph{स्वयंसिद्धकीय पद्धत} ही आधारित असते: एखाद्या
संरचनेच्या संग्रहाचे उपपत्तीच्या स्वयंसिद्धकांचा, म्हणजे
वाक्यांच्या संचाचा, वापर करून वर्णन करणे. स्वयंसिद्धकांपासून
प्रथम-क्रम तर्कशास्त्रात तार्किकरीत्या निष्पन्न होणारी नेमकी तीच
वाक्ये स्वयंसिद्धकांच्या सर्व प्रतिमानांत सत्य असतात, या
निरीक्षणामुळे हे शक्य होते.

अतिशय सोपे उदाहरण म्हणून पूर्वक्रमांचा विचार करा. एखाद्या संच~$A$ वरील
परावर्ती आणि संक्रमक असा संबंध~$R$ म्हणजे पूर्वक्रम होय. द्विस्थानी
संबंध $R \subseteq A \times A$ असलेला संच~$A$ म्हणजे एकच द्विस्थानी
संबंधचिन्ह~$\Obj P$ असलेल्या प्रथम-क्रम भाषेला संरचना देण्यासाठी
नेमके जे आवश्यक आहे ते: आपण $\Domain{M} = A$ आणि
$\Assign{\Obj P}{M} = R$ ठेवू. $R$ हा पूर्वक्रम असल्यामुळे तो परावर्ती
आणि संक्रमक आहे आणि हे सांगणारा प्रथम-क्रम तर्कशास्त्रातील
वाक्यांचा संच~$\Gamma$ आपल्याला मिळू शकतो:
\begin{align*}
  & \lforall[\Obj v_0][\Atom{\Obj P}{\Obj v_0,\Obj v_0}]\\
  & \lforall[\Obj v_0][\lforall[\Obj v_1][\lforall[\Obj v_2][((\Atom{\Obj P}{\Obj v_0,\Obj v_1} \land \Atom{\Obj P}{\Obj v_1,\Obj v_2}) \lif \Atom{\Obj P}{\Obj v_0,\Obj v_2})]]]
\end{align*}
ही वाक्ये म्हणजे फक्त ``कोणत्याही~$x$ साठी $Rxx$'' ($R$ हा
परावर्ती आहे) आणि ``जेव्हा $Rxy$ आणि $Ryz$, तेव्हा $Rxz$ ही'' ($R$ हा
संक्रमक आहे) यांची चिन्हांकने आहेत. संरचना~$\Struct{M}$ हे या
दोन वाक्ये~$\Gamma$ चे प्रतिमान तेव्हा आणि केवळ तेव्हाच असते,
जेव्हा $R$ (म्हणजे $\Assign{\Obj P}{M}$) हा~$A$ (म्हणजे
$\Domain{M}$) वरील पूर्वक्रम असतो. दुसऱ्या शब्दांत, $\Gamma$ ची प्रतिमाने
म्हणजे नेमके पूर्वक्रम होत. फक्त~$\Obj P$ हे विधेय असलेल्या
प्रथम-क्रम भाषेत व्यक्त करता येणारा सर्व पूर्वक्रमांचा कोणताही गुणधर्म
(वरील परावर्तिता आणि संक्रमणता यांसारखा) हा~$\Gamma$ मधील दोन
वाक्यांपासून तार्किकरीत्या निष्पन्न होतो आणि उलटही तसेच होते.
म्हणून~$\Gamma$ च्या प्रतिमानांविषयी आपण जे काही सिद्ध करू शकतो, ते सर्व
पूर्वक्रमांविषयी सिद्ध केलेले असते.

कोणत्याही विशिष्ट उपपत्तीसाठी आणि प्रतिमानांच्या वर्गासाठी (जसे
$\Gamma$ आणि सर्व पूर्वक्रम) संबंधित प्रथम-क्रम भाषेत काय व्यक्त करता
येते आणि काय व्यक्त करता येत नाही याविषयी रोचक प्रश्न असतील. सर्व भाषा
आणि प्रतिमानवर्गांसाठी रोचक असणारे संरचनेचे काही गुणधर्म आहेत;
विशेषतः प्रांताच्या आकारमानाशी संबंधित गुणधर्म. उदाहरणार्थ,
कोणत्याही~$n \in \PosInt$ साठी प्रांत मध्ये नेमके
$n$~घटक आहेत असे नेहमी व्यक्त करता येते. अनंतपणे अनेक
वाक्यांचा संच वापरून प्रांत अनंत आहे असेही व्यक्त करता येते.
पण प्रांत सांत आहे किंवा प्रांत अगणनीय आहे असे व्यक्त करता
येत नाही. प्रथम-क्रम भाषांच्या मर्यादांविषयीचे हे परिणाम संहतता आणि
लोव्हेनहाइम--स्कोलेम प्रमेयांचे परिणाम आहेत.












\section{निर्दोषता आणि संपूर्णता}\label{fol:int:scp:sec}

प्रथम-क्रम तर्कशास्त्रासाठीच्या निष्पत्ती-पद्धतींचाही आपण परिचय करून
घेऊ. तर्कशास्त्रज्ञांनी अनेक निष्पत्ती-पद्धती विकसित केल्या आहेत,
पण त्या सर्व वाक्ये दरम्यान तोच निष्पन्नता संबंध परिभाषित
करतात. विशिष्ट आणि काटेकोरपणे परिभाषित प्रकारची निष्पत्ती असेल,
तर~$\Gamma$ हे~$\varphi$ \emph{निष्पन्न} करते, $\Gamma \Proves \varphi$, असे
आपण म्हणतो. निष्पत्त्या म्हणजे नेहमी चिन्हांच्या सांत मांडण्या
असतात---कदाचित वाक्यांची यादी किंवा त्याहून गुंतागुंतीची काही
संरचना. एखाद्या संच~$\Gamma$ पासून वाक्य तार्किकरीत्या निष्पन्न
होते का हे ठरवण्याचे साधन पुरवणे हे निष्पत्ती-पद्धतींचे उद्दिष्ट
आहे. ते उद्दिष्ट साध्य करण्यासाठी $\Gamma \Entails \varphi$ तेव्हा आणि केवळ
तेव्हाच $\Gamma \Proves \varphi$ असले पाहिजे.

$\Gamma \Proves \varphi$ पण $\Gamma \Entails \varphi$ नसेल, तर आपली
निष्पत्ती-पद्धत अतिप्रबल असेल, म्हणजे ती गरजेपेक्षा जास्त सिद्ध
करेल. $\Gamma \Proves \varphi$ असेल तर $\Gamma \Entails \varphi$ असते या
गुणधर्माला \emph{निर्दोषता} म्हणतात आणि कोणत्याही चांगल्या
निष्पत्ती-पद्धतीसाठी ही किमान आवश्यकता आहे. दुसरीकडे,
$\Gamma \Entails \varphi$ पण $\Gamma \Proves \varphi$ नसेल, तर आपली
निष्पत्ती-पद्धत अतिदुर्बल असेल, म्हणजे ती पुरेसे सिद्ध करणार नाही.
$\Gamma \Entails \varphi$ असेल तर $\Gamma \Proves \varphi$ असते या गुणधर्माला
\emph{संपूर्णता} म्हणतात. निर्दोषता सिद्ध करणे सहसा तुलनेने सोपे असते
(विगमनाधारित रीतीने परिभाषित केलेल्या निष्पत्त्यांच्या संरचनेवर
विगमन करून). संपूर्णता सिद्ध करणे अधिक कठीण असते.

निर्दोषता आणि संपूर्णता यांचे अनेक महत्त्वाचे परिणाम आहेत.
वाक्यांचा एखादा संच~$\Gamma$ हा व्याघात (जसे
$\varphi \land \lnot \varphi$) निष्पन्न करत असेल, तर त्याला \emph{विसंगत}
म्हणतात. विसंगत~$\Gamma$ ला कोणतेही प्रतिमान असू शकत नाही; तो
अपूर्ततायोग्य असतो. संपूर्णतेवरून याचा उलट निष्कर्ष मिळतो: विसंगत
नसलेल्या---किंवा, आपण म्हणू त्याप्रमाणे, \emph{सुसंगत}---प्रत्येक
$\Gamma$ ला प्रतिमान असते. खरे तर, हे संपूर्णतेशी सममूल्य आहे आणि
संपूर्णतेचे हेच रूप आपण प्रत्यक्षात सिद्ध करू. हा सखोल आणि कदाचित
आश्चर्यकारक परिणाम आहे: $\Gamma$ पासून $\varphi \land \lnot \varphi$ सिद्ध करता
येत नाही एवढ्यामुळेच~$\Gamma$ वर्णन करते तसे संरचना अस्तित्वात
असण्याची हमी मिळते. म्हणून संपूर्णता पुढील प्रश्नाचे उत्तर देते:
वाक्यांच्या कोणत्या संचांना प्रतिमाने असतात? उत्तर: सर्व आणि फक्त
सुसंगत संचांना.

निर्दोषता आणि संपूर्णता प्रमेयांचे दोन महत्त्वाचे परिणाम आहेत: संहतता
प्रमेय आणि लोव्हेनहाइम--स्कोलेम प्रमेय. प्रतिमानांच्या उपपत्तीतील हे
महत्त्वाचे परिणाम आहेत आणि अनेक रोचक निष्कर्ष प्रस्थापित करण्यासाठी
त्यांचा उपयोग करता येतो. आपण आधीच दोन निष्कर्षांचा उल्लेख केला आहे:
प्रथम-क्रम तर्कशास्त्रात संरचनेचा प्रांत सांत आहे किंवा तो
अगणनीय आहे असे व्यक्त करता येत नाही.

ऐतिहासिकदृष्ट्या हे सर्व---प्रथम-क्रम तर्कशास्त्राची विन्यासमीमांसा आणि
चिन्हार्थमीमांसा कशी परिभाषित करायची, चांगल्या निष्पत्ती-पद्धती
कशा परिभाषित करायच्या, त्या निर्दोष आणि संपूर्ण आहेत हे कसे सिद्ध
करायचे, प्रथम-क्रम भाषांत काय व्यक्त करता येते आणि काय व्यक्त करता येत
नाही हे स्पष्ट करणे---शोधून योग्य रीतीने ठरवण्यास बराच काळ लागला. हे
कसे करायचे हे आता माहीत असले, तरी सर्व तपशील समजून घेणे अजूनही
गोंधळात टाकणारे आणि कंटाळवाणे ठरू शकते. पण ते महत्त्वाचेही आहे, कारण
प्रथम-क्रम तर्कशास्त्राच्या आकारिक भाषेसाठी येथे विकसित केलेल्या पद्धती
तर्कशास्त्र, संगणकशास्त्र आणि भाषाविज्ञान यांत सर्वत्र वापरल्या जातात.
म्हणून सर्व तपशीलांचा अभ्यास करण्याचे दीर्घकाळात फळ मिळते.



